\documentclass{amsart}
% For dealing with references we use the comment environment
\usepackage{verbatim}
\newenvironment{reference}{\comment}{\endcomment}
%\newenvironment{reference}{}{}
\newenvironment{slogan}{\comment}{\endcomment}
\newenvironment{history}{\comment}{\endcomment}

% For commutative diagrams we use Xy-pic
\usepackage[all]{xy}

% We use 2cell for 2-commutative diagrams.
\xyoption{2cell}
\UseAllTwocells

% We use multicol for the list of chapters between chapters
\usepackage{multicol}

% This is generally recommended for better output
\usepackage{lmodern}
\usepackage[T1]{fontenc}

% For cross-file-references

% Package for hypertext links:
\usepackage{hyperref}

% For any local file, say "hello.tex" you want to link to please

% Theorem environments.
%
\theoremstyle{plain}
\newtheorem{theorem}[subsection]{Theorem}
\newtheorem{proposition}[subsection]{Proposition}
\newtheorem{lemma}[subsection]{Lemma}

\theoremstyle{definition}
\newtheorem{definition}[subsection]{Definition}
\newtheorem{example}[subsection]{Example}
\newtheorem{exercise}[subsection]{Exercise}
\newtheorem{situation}[subsection]{Situation}

\theoremstyle{remark}
\newtheorem{remark}[subsection]{Remark}
\newtheorem{remarks}[subsection]{Remarks}

\numberwithin{equation}{subsection}

% Macros
%
\def\lim{\mathop{\mathrm{lim}}\nolimits}
\def\colim{\mathop{\mathrm{colim}}\nolimits}
\def\Spec{\mathop{\mathrm{Spec}}}
\def\Hom{\mathop{\mathrm{Hom}}\nolimits}
\def\Ext{\mathop{\mathrm{Ext}}\nolimits}
\def\SheafHom{\mathop{\mathcal{H}\!\mathit{om}}\nolimits}
\def\SheafExt{\mathop{\mathcal{E}\!\mathit{xt}}\nolimits}
\def\Sch{\mathit{Sch}}
\def\Mor{\mathop{\mathrm{Mor}}\nolimits}
\def\Ob{\mathop{\mathrm{Ob}}\nolimits}
\def\Sh{\mathop{\mathit{Sh}}\nolimits}
\def\NL{\mathop{N\!L}\nolimits}
\def\CH{\mathop{\mathrm{CH}}\nolimits}
\def\proetale{{pro\text{-}\acute{e}tale}}
\def\etale{{\acute{e}tale}}
\def\QCoh{\mathit{QCoh}}
\def\Ker{\mathop{\mathrm{Ker}}}
\def\Im{\mathop{\mathrm{Im}}}
\def\Coker{\mathop{\mathrm{Coker}}}
\def\Coim{\mathop{\mathrm{Coim}}}

% Boxtimes
%
\DeclareMathSymbol{\boxtimes}{\mathbin}{AMSa}{"02}

%
% Macros for moduli stacks/spaces
%
\def\QCohstack{\mathcal{QC}\!\mathit{oh}}
\def\Cohstack{\mathcal{C}\!\mathit{oh}}
\def\Spacesstack{\mathcal{S}\!\mathit{paces}}
\def\Quotfunctor{\mathrm{Quot}}
\def\Hilbfunctor{\mathrm{Hilb}}
\def\Curvesstack{\mathcal{C}\!\mathit{urves}}
\def\Polarizedstack{\mathcal{P}\!\mathit{olarized}}
\def\Complexesstack{\mathcal{C}\!\mathit{omplexes}}
% \Pic is the operator that assigns to X its picard group, usage \Pic(X)
% \Picardstack_{X/B} denotes the Picard stack of X over B
% \Picardfunctor_{X/B} denotes the Picard functor of X over B
\def\Pic{\mathop{\mathrm{Pic}}\nolimits}
\def\Picardstack{\mathcal{P}\!\mathit{ic}}
\def\Picardfunctor{\mathrm{Pic}}
\def\Deformationcategory{\mathcal{D}\!\mathit{ef}}

\setlength{\emergencystretch}{3em}
\begin{document}
\title[Stacks Project, Tag 0DQF]{403 --- Versal formal objects after a finite extension of the residue field}
\maketitle
\noindent Copyright (C) 2005--2025 Johan de Jong. Stacks Project authors. Source: \href{https://stacks.math.columbia.edu/tag/0DQF}{Tag 0DQF}.

\noindent Editorial difficulty note (not author text). Difficulty H2: The proof constructs a formally smooth completed polynomial extension realizing the new residue field. It then defines a natural transformation into the residue-change functor and verifies its lifting property against arbitrary Artinian surjections using continuous formal smoothness. This establishes compatibility of versality with the nontrivial change-of-residue-field construction..

\noindent Editorial provenance note (not author text). History: excerpt from revision \texttt{a0444\allowbreak 6e57e\allowbreak c1fbc\allowbreak 252a8\allowbreak 71afc\allowbreak ec775\allowbreak 2fb28\allowbreak 07b14}, formal-defos.tex, lines 7475--7536. Reference presentation was adapted; original-author context and core support blocks were retained or added. The mathematical wording is unchanged. Licensed under GNU FDL 1.2 or later; no invariant sections or cover texts. See ../licenses/COPYING. Context blocks reproduce author definitions and supporting results; they are not additional counted problems.

\begin{definition}
\label{definition-CLambda}
Let $\Lambda$ be a Noetherian ring and let $\Lambda \to k$ be a finite
ring map where $k$ is a field. We define {\it $\mathcal{C}_\Lambda$} to be
the category with
\begin{enumerate}
\item objects are pairs $(A, \varphi)$ where $A$ is an Artinian local
$\Lambda$-algebra and where $\varphi : A/\mathfrak m_A \to k$ is a
$\Lambda$-algebra isomorphism, and
\item morphisms $f : (B, \psi) \to (A, \varphi)$ are local $\Lambda$-algebra
homomorphisms such that $\varphi \circ (f \bmod \mathfrak m) = \psi$.
\end{enumerate}
We say we are in the {\it classical case} if $\Lambda$ is a Noetherian
complete local ring and $k$ is its residue field.
\end{definition}

\section{Notation and Conventions}
\label{section-notations-conventions}

\noindent
A ring is commutative with $1$. The  maximal ideal of a local ring $A$
is denoted by $\mathfrak{m}_A$. The set of positive integers is denoted
by $\mathbf{N} = \{1, 2, 3, \ldots\}$. If $U$ is an object of a
category $\mathcal{C}$, we denote by $\underline{U}$
the functor
$\Mor_\mathcal{C}(U, -): \mathcal{C} \to \textit{Sets}$, see
Remarks \href{https://stacks.math.columbia.edu/tag/06GK}{remarks cofibered groupoids (Tag 06GK)} (\href{https://stacks.math.columbia.edu/tag/06GQ}{item definition yoneda (Tag 06GQ)}).
Warning: this may conflict with the notation in other chapters where we
sometimes use $\underline{U}$ to denote $h_U(-) = \Mor_\mathcal{C}(-, U)$.

\medskip\noindent
Throughout this chapter $\Lambda$ is a Noetherian ring and
$\Lambda \to k$ is a finite ring map from $\Lambda$ to a field.
The kernel of this map is denoted $\mathfrak m_\Lambda$ and the
image $k' \subset k$. It turns out that $\mathfrak m_\Lambda$ is
a maximal ideal, $k' = \Lambda/\mathfrak m_\Lambda$ is a field, and
the extension $k/k'$ is finite. See discussion surrounding
(\href{https://stacks.math.columbia.edu/tag/06S2}{equation k prime (Tag 06S2)}).




\begin{definition}
\label{definition-completion-CLambda}
Let $\Lambda$ be a Noetherian ring and let $\Lambda \to k$ be a finite
ring map where $k$ is a field. We define {\it $\widehat{\mathcal{C}}_\Lambda$}
to be the category with
\begin{enumerate}
\item objects are pairs $(R, \varphi)$ where $R$ is a Noetherian complete
local $\Lambda$-algebra and where $\varphi : R/\mathfrak m_R \to k$ is a
$\Lambda$-algebra isomorphism, and
\item morphisms $f : (S, \psi) \to (R, \varphi)$ are local $\Lambda$-algebra
homomorphisms such that $\varphi \circ (f \bmod \mathfrak m) = \psi$.
\end{enumerate}
\end{definition}

\begin{definition}
\label{definition-category-cofibred-groupoids}
Let $\mathcal{C}$ be a category.  A {\it category cofibered in groupoids over
$\mathcal{C}$} is a category $\mathcal{F}$ equipped with a functor
$p: \mathcal{F} \to \mathcal{C}$ such that $\mathcal{F}^{opp}$ is a category
fibered in groupoids over $\mathcal{C}^{opp}$ via
$p^{opp}: \mathcal{F}^{opp} \to \mathcal{C}^{opp}$.
\end{definition}

\begin{definition}
\label{definition-predeformation-category}
A {\it predeformation category} $\mathcal{F}$ is a category cofibered
in groupoids over $\mathcal{C}_\Lambda$ such that $\mathcal{F}(k)$ is
equivalent to a category with a single object and a single morphism,
i.e., $\mathcal{F}(k)$ contains at least one object and there is a
unique morphism between any two objects. A {\it morphism of predeformation
categories} is a morphism of categories cofibered in groupoids over
$\mathcal{C}_\Lambda$.
\end{definition}

\begin{definition}
\label{definition-formal-objects}
Let $\mathcal{F}$ be a category cofibered in groupoids over
$\mathcal{C}_\Lambda$. The {\it category $\widehat{\mathcal{F}}$ of formal
objects of  $\mathcal{F}$} is the category with the following objects and
morphisms.
\begin{enumerate}
\item A {\it formal object $\xi = (R, \xi_n, f_n)$ of $\mathcal{F}$}
consists of an object $R$ of $\widehat{\mathcal{C}}_\Lambda$, and a collection
indexed by $n \in \mathbf{N}$ of objects $\xi_n$ of
$\mathcal{F}(R/\mathfrak m_R^n)$ and morphisms
$f_n : \xi_{n + 1} \to \xi_n$ lying over the projection
$R/\mathfrak m_R^{n + 1} \to R/\mathfrak m_R^n$.
\item Let $\xi = (R, \xi_n, f_n)$ and $\eta = (S, \eta_n, g_n)$ be
formal objects of $\mathcal{F}$.  A {\it morphism $a : \xi \to \eta$ of
formal objects} consists of a map $a_0 : R \to S$ in
$\widehat{\mathcal{C}}_\Lambda$ and a collection $a_n : \xi_n \to \eta_n$
of morphisms of $\mathcal{F}$ lying over
$R/\mathfrak m_R^n \to S/\mathfrak m_S^n$,
such that for every $n$ the diagram
$$
\xymatrix{
\xi_{n + 1} \ar[r]_{f_n} \ar[d]_{a_{n + 1}} & \xi_n \ar[d]^{a_n} \\
\eta_{n + 1} \ar[r]^{g_n} & \eta_n
}
$$
commutes.
\end{enumerate}
\end{definition}

\begin{definition}
\label{definition-versal}
Let $\mathcal{F}$ be a category cofibered in groupoids.  Let $\xi$ be a formal
object of $\mathcal{F}$ lying over $R \in \Ob(\widehat{\mathcal{C}}_\Lambda)$.
We say $\xi$ is {\it versal} if the corresponding morphism
$\underline{\xi}: \underline{R}|_{\mathcal{C}_\Lambda} \to \mathcal{F}$
of Remark \href{https://stacks.math.columbia.edu/tag/06HC}{remark formal objects yoneda (Tag 06HC)} is smooth.
\end{definition}

\begin{definition}
\label{definition-S1-S2}
Let $\mathcal{F}$ be a category cofibered in groupoids over $\mathcal
C_\Lambda$. We define {\it conditions (S1) and (S2)}
on $\mathcal{F}$ as follows:
\begin{enumerate}
\item[(S1)] Every diagram in $\mathcal{F}$
$$
\vcenter{
\xymatrix{
           & x_2 \ar[d] \\
x_1 \ar[r] & x
}
}
\quad\text{lying over}\quad
\vcenter{
\xymatrix{
           & A_2 \ar[d] \\
A_1 \ar[r] & A
}
}
$$
in $\mathcal{C}_\Lambda$ with $A_2 \to A$ surjective can be completed
to a commutative diagram
$$
\vcenter{
\xymatrix{
y \ar[r] \ar[d] & x_2 \ar[d] \\
x_1 \ar[r]      & x
}
}
\quad\text{lying over}\quad
\vcenter{
\xymatrix{
A_1 \times_A A_2 \ar[r] \ar[d] & A_2 \ar[d] \\
A_1 \ar[r]      & A.
}
}
$$
\item[(S2)]
The condition of (S1) holds for diagrams in $\mathcal{F}$ lying over
a diagram in $\mathcal{C}_\Lambda$ of the form
$$
\xymatrix{
          & k[\epsilon] \ar[d] \\
A  \ar[r] & k.
}
$$
Moreover, if we have two commutative diagrams in $\mathcal{F}$
$$
\vcenter{
\xymatrix{
y \ar[r]_c \ar[d]_a & x_\epsilon \ar[d]^e \\
x \ar[r]^d          & x_0
}
}
\quad\text{and}\quad
\vcenter{
\xymatrix{
y' \ar[r]_{c'} \ar[d]_{a'} & x_\epsilon \ar[d]^e \\
x \ar[r]^d                 & x_0
}
}
\quad\text{lying over}\quad
\vcenter{
\xymatrix{
A \times_k k[\epsilon] \ar[r] \ar[d] & k[\epsilon] \ar[d] \\
A  \ar[r] & k
}
}
$$
then there exists a morphism $b : y \to y'$ in
$\mathcal{F}(A \times_k k[\epsilon])$ such that $a = a' \circ b$.
\end{enumerate}
\end{definition}

\begin{situation}
\label{situation-change-of-fields}
Let $\Lambda$ be a Noetherian ring and let $\Lambda \to k \to l$ be a finite
ring maps where $k$ and $l$ are fields. Thus $l/k$ is a finite
extensions of fields. A typical object of $\mathcal{C}_{\Lambda, l}$ will be
denoted $B$ and a typical object of $\mathcal{C}_{\Lambda, k}$ will be
denoted $A$. We define
\begin{equation}
\label{equation-comparison}
\mathcal{C}_{\Lambda, l} \longrightarrow \mathcal{C}_{\Lambda, k},
\quad
B \longmapsto B \times_l k
\end{equation}
Given a category cofibred in groupoids
$p : \mathcal{F} \to \mathcal{C}_{\Lambda, k}$ we obtain
an associated category cofibred in groupoids
$$
p_{l/k} : \mathcal{F}_{l/k} \longrightarrow \mathcal{C}_{\Lambda, l}
$$
by setting $\mathcal{F}_{l/k}(B) = \mathcal{F}(B \times_l k)$.
\end{situation}

\noindent
Explicitly, $p: \mathcal{F} \to \mathcal{C}$ is cofibered in groupoids if
the following two conditions hold:
\begin{enumerate}
\item For every morphism $f: U \to V$ in $\mathcal{C}$ and every object
$x$ lying over $U$, there is a morphism $x \to y$ of $\mathcal{F}$ lying
over $f$.
\item For every pair of morphisms $a: x \to y$ and $b: x \to z$
of $\mathcal{F}$ and any morphism $f: p(y) \to p(z)$ such that $p(b) = f
\circ p(a)$, there exists a unique morphism $c: y \to z$ of $\mathcal
F$ lying over $f$ such that $b = c \circ a$.
\end{enumerate}


\begin{lemma}
\label{lemma-change-of-field-versal-ring}
With notation and assumptions as in Situation \ref{situation-change-of-fields}.
Let $\xi$ be a versal formal object for $\mathcal{F}$ lying over
$R \in \Ob(\widehat{\mathcal{C}}_{\Lambda, k})$. Then there exist
\begin{enumerate}
\item an $S \in \Ob(\widehat{\mathcal{C}}_{\Lambda, l})$
and a local $\Lambda$-algebra homomorphism $R \to S$ which is
formally smooth in the $\mathfrak m_S$-adic topology and induces
the given field extension $l/k$ on residue fields, and
\item a versal formal object of $\mathcal{F}_{l/k}$ lying over $S$.
\end{enumerate}
\end{lemma}

\begin{proof}
Construction of $S$. Choose a surjection $R[x_1, \ldots, x_n] \to l$
of $R$-algebras. The kernel is a maximal ideal $\mathfrak m$.
Set $S$ equal to the $\mathfrak m$-adic completion of the Noetherian ring
$R[x_1, \ldots, x_n]$. Then $S$ is in $\widehat{\mathcal{C}}_{\Lambda, l}$
by Algebra, Lemma \href{https://stacks.math.columbia.edu/tag/0316}{lemma completion Noetherian Noetherian (Tag 0316)}.
The map $R \to S$ is formally smooth in the $\mathfrak m_S$-adic topology by
More on Algebra, Lemmas
\href{https://stacks.math.columbia.edu/tag/07EC}{lemma formally smooth (Tag 07EC)} and
\href{https://stacks.math.columbia.edu/tag/07ED}{lemma formally smooth completion (Tag 07ED)}
and the fact that $R \to R[x_1, \ldots, x_n]$ is formally smooth.
(Compare with the proof Lemma \href{https://stacks.math.columbia.edu/tag/06SM}{lemma exists smooth (Tag 06SM)}.)

\medskip\noindent
Since $\xi$ is versal, the transformation
$\underline{\xi} : \underline{R}|_{\mathcal{C}_{\Lambda, k}} \to \mathcal{F}$
is smooth. By Lemma \href{https://stacks.math.columbia.edu/tag/07WD}{lemma change of fields smooth (Tag 07WD)} the induced map
$$
(\underline{R}|_{\mathcal{C}_{\Lambda, k}})_{l/k}
\longrightarrow
\mathcal{F}_{l/k}
$$
is smooth. Thus it suffices to construct a smooth morphism
$\underline{S}|_{\mathcal{C}_{\Lambda, l}} \to
(\underline{R}|_{\mathcal{C}_{\Lambda, k}})_{l/k}$.
To give such a map means for every object $B$ of $\mathcal{C}_{\Lambda, l}$
a map of sets
$$
\Mor_{\widehat{\mathcal{C}}_{\Lambda, l}}(S, B)
\longrightarrow
\Mor_{\widehat{\mathcal{C}}_{\Lambda, k}}(R, B \times_l k)
$$
functorial in $B$. Given an element $\varphi : S \to B$
on the left hand side we send it to the composition $R \to S \to B$
whose image is contained in the sub $\Lambda$-algebra $B \times_l k$.
Smoothness of the map means that given a surjection
$B' \to B$ and a commutative diagram
$$
\xymatrix{
S \ar[r] & B \ar@{=}[r] & B \\
R \ar[u] \ar[r] & B' \times_l k \ar[u] \ar[r] & B' \ar[u]
}
$$
we have to find a ring map $S \to B'$ fitting into the outer rectangle.
The existence of this map is guaranteed as we chose $R \to S$ to be
formally smooth in the $\mathfrak m_S$-adic topology, see
More on Algebra, Lemma \href{https://stacks.math.columbia.edu/tag/07NJ}{lemma lift continuous (Tag 07NJ)}.
\end{proof}
\end{document}
